\documentclass[11pt,a4paper]{article}
\usepackage[utf8]{inputenc}
\usepackage[T1]{fontenc}
\usepackage{amsmath,amssymb,amsthm}
\usepackage[margin=2.6cm]{geometry}
\usepackage{hyperref}
\usepackage{booktabs}
\usepackage{xcolor}
\usepackage{listings}
\lstset{basicstyle=\ttfamily\small,breaklines=true,columns=fullflexible,keepspaces=true,escapechar=@}

\newtheorem{theorem}{Theorem}[section]
\newtheorem{lemma}[theorem]{Lemma}
\newtheorem{proposition}[theorem]{Proposition}
\newtheorem{corollary}[theorem]{Corollary}
\theoremstyle{definition}
\newtheorem{definition}[theorem]{Definition}
\theoremstyle{remark}
\newtheorem{remark}[theorem]{Remark}

\newcommand{\N}{\mathbb{N}}
\newcommand{\R}{\mathbb{R}}
\newcommand{\C}{\mathbb{C}}
\newcommand{\Z}{\mathbb{Z}}
\DeclareMathOperator{\tr}{tr}
\DeclareMathOperator{\charp}{charpoly}
\DeclareMathOperator{\detm}{det}

\title{Conjecture 00000003718:\\ the matching polynomial of a tree has only real roots\\
\large A proof, with a Lean~4 / Mathlib formalization}
\author{Zhiyuan Zhang (GitHub: \texttt{LibertychaserUS})}
\date{27 October 2026}

\begin{document}
\maketitle

\begin{abstract}
We prove conjecture 00000003718 of \emph{The Last Math Competition}: for every finite tree $T$
the matching polynomial $\mu(T;x)=\sum_{k\ge 0}(-1)^k m_k(T)\,x^{n-2k}$ has only real roots, and
the distribution of these roots is determined by the walk counts of $T$. The proof goes through
the identity $\mu(F;x)=\det(xI-A(F))$ for every forest $F$, established by a leaf-deletion
induction on both sides, followed by the spectral theorem for real symmetric matrices; the moment
identity $\sum_i\theta_i^{\,k}=\tr A(T)^k=\#\{\text{closed walks of length }k\}$ identifies the
root distribution with the moment transform of the closed-walk counts. Every statement is
formalized and machine-checked in Lean~4 (toolchain \texttt{v4.33.1}, Mathlib \texttt{v4.33.1});
the project compiles with no \texttt{sorry}, no \texttt{native\_decide} and no axioms beyond
\texttt{propext}, \texttt{Classical.choice} and \texttt{Quot.sound}.
\end{abstract}

\section{The conjecture as filed}

\begin{quote}\small
\textbf{English.} Definition: Matching polynomials of trees. Conjecture: The real roots of
matching polynomials: the roots of the matching polynomial of trees are all real, with the
distribution of real roots the transform of path counts. (tree matching real root path count
transform distribution)

\emph{(The Chinese version of the statement in \texttt{conjectures/00000003718.md} is a literal
translation of the English one.)}
\end{quote}

\subsection{Reading of the statement}

Throughout, a \emph{graph} is a finite simple graph $G=(V,E)$ with $n=|V|$ vertices; a
\emph{tree} is a connected acyclic graph, a \emph{forest} an acyclic graph.

\begin{definition}[matching polynomial]
A \emph{$k$-matching} of $G$ is a set of $k$ pairwise vertex-disjoint edges; write $m_k(G)$ for
the number of $k$-matchings ($m_0=1$). The \emph{matching polynomial} of $G$ is
\[
  \mu(G;x) \;=\; \sum_{k\ge 0} (-1)^k\, m_k(G)\, x^{\,n-2k}
  \;=\; \sum_{M\ \text{matching of }G} (-1)^{|M|} x^{\,n-2|M|}\in\Z[x].
\]
\end{definition}
This is the standard definition (Heilmann--Lieb 1972; Farrell 1979; Godsil--Gutman 1981), and it is
exactly the one formalized in \texttt{SimpleGraph.matchingPolynomial}; the equality of the two
displayed forms is the Lean theorem \texttt{matchingPolynomial\_eq\_sum\_numMatchings}.

The conjecture has two clauses.

\begin{enumerate}
\item[(C1)] \emph{``the roots of the matching polynomial of trees are all real''}: for every tree
$T$, every (complex) root of $\mu(T;x)$ is real; equivalently $\mu(T;x)$ splits into $n$ linear
factors over $\R$.
\item[(C2)] \emph{``with the distribution of real roots the transform of path counts''}: the
multiset of roots $\{\theta_1,\dots,\theta_n\}$ is determined by counting walks in $T$. We make
this precise through its moments: for every $k\ge 0$,
\[
  \sum_{i=1}^n \theta_i^{\,k} \;=\; \#\{\text{closed walks of length $k$ in }T\}
  \;=\; \sum_{v\in V}\#\{\text{walks }v\to v\text{ of length }k\}.
\]
The power sums of a degree-$n$ monic polynomial determine it (Newton's identities), so the
root distribution is literally the moment transform of the closed-walk counting sequence; in a
tree every closed walk is a back-and-forth traversal of a path, which is the ``path count''
phrasing of the statement.
\end{enumerate}
We prove both clauses. The precise formal statement is Theorem~\ref{thm:main} below, whose Lean
counterpart is \texttt{TLMC3718.conjecture\_00000003718}.

\begin{remark}
Clause (C1) is a classical theorem: Heilmann and Lieb (1972) proved real-rootedness of
$\mu(G;x)$ for \emph{every} graph, and Godsil and Gutman (1981) observed that for forests
$\mu(F;x)=\det(xI-A(F))$, which gives the tree case directly. We give a complete self-contained
proof of the tree case along the second route, and we formalize it; we are not aware of a
previous formalization of these results in Mathlib.
\end{remark}

\section{Statement of the results}

Let $A=A(G)\in\{0,1\}^{V\times V}$ be the adjacency matrix and $\charp(A)=\det(xI-A)\in\Z[x]$.
For $S\subseteq V$ let $G[S]$ be the induced subgraph, $A_S$ its adjacency matrix (the principal
submatrix of $A$ indexed by $S$) and $\mu_S(x)=\mu(G[S];x)$; a matching of $G[S]$ is a set of
edges of $G$ with both end points in $S$, pairwise disjoint.

\begin{theorem}[forests: $\mu=\charp$]\label{thm:A}
Let $G$ be a forest. For every $S\subseteq V$, $\ \mu(G[S];x)=\det(xI-A_S)$. In particular
$\mu(G;x)=\charp(A(G))$.
\end{theorem}

\begin{theorem}[main theorem; conjecture 00000003718]\label{thm:main}
Let $T$ be a tree on $n$ vertices with roots $\theta_1,\dots,\theta_n$ of $\mu(T;x)$ in $\C$
(with multiplicity). Then
\begin{enumerate}
\item[(a)] $\mu(T;x)$ splits over $\R$: $\mu(T;x)=\prod_{i=1}^n (x-\theta_i)$ with all
$\theta_i\in\R$; in particular $\mu(T;x)$ has exactly $n$ real roots counted with multiplicity,
and every complex root $z$ of $\mu(T;x)$ satisfies $\operatorname{Im} z=0$;
\item[(b)] for every $k\in\N$,
$\ \sum_{i=1}^n\theta_i^{\,k}=\sum_{v\in V}\#\{\text{walks }v\to v\text{ in $T$ of length }k\}$.
\end{enumerate}
\end{theorem}

\section{Proofs}

\subsection{Leaves of forests}

\begin{lemma}\label{lem:leaf}
Let $G$ be a forest and $S\subseteq V$. If $G$ has an edge with both end points in $S$, then there
is $v\in S$ having exactly one neighbour in $S$.
\end{lemma}
\begin{proof}
$G[S]$ is again a forest, and it has an edge, so it has a path of length $\ge1$. Let $P$ be a
path of maximal length in $G[S]$, with first vertex $v$ and second vertex $u$. If $w\in S$ is a
neighbour of $v$ with $w\notin P$, then prepending $w$ to $P$ gives a longer path, contradicting
maximality; so every neighbour of $v$ in $S$ lies on $P$. In an acyclic graph the only vertex of
$P$ adjacent to its first vertex $v$ is the second vertex $u$ (otherwise the subpath of $P$ from
$v$ to $w$ together with the edge $wv$ is a cycle). Hence $u$ is the unique neighbour of $v$ in $S$.
\end{proof}

\subsection{The leaf recurrence for the matching polynomial}

\begin{lemma}\label{lem:mu}
Let $S\subseteq V$, $v\in S$, and suppose $u\in S$ is the unique neighbour of $v$ in $S$. Then
\[
  \mu_S(x)\;=\;x\,\mu_{S\setminus\{v\}}(x)\;-\;\mu_{S\setminus\{v,u\}}(x).
\]
\end{lemma}
\begin{proof}
Split the matchings $M$ of $G[S]$ according to whether $vu\in M$.
If $vu\notin M$ then no edge of $M$ contains $v$ (such an edge would be $vw$ with $w\in S$ a
neighbour of $v$, hence $w=u$), so $M$ is a matching of $G[S\setminus\{v\}]$; conversely every
matching of $G[S\setminus\{v\}]$ is a matching of $G[S]$ not containing $vu$. Since
$|S|=|S\setminus\{v\}|+1$ and $2|M|\le|S\setminus\{v\}|$ (a matching inside a set $S'$ covers
$2|M|$ distinct vertices of $S'$), these matchings contribute
$\sum_M(-1)^{|M|}x^{|S\setminus\{v\}|+1-2|M|}=x\,\mu_{S\setminus\{v\}}(x)$.
If $vu\in M$ then $M'=M\setminus\{vu\}$ is a matching of $G[S\setminus\{v,u\}]$ (its edges are
disjoint from $vu$), and $M\mapsto M'$ is a bijection onto the matchings of
$G[S\setminus\{v,u\}]$ with inverse $M'\mapsto M'\cup\{vu\}$; as $|S|=|S\setminus\{v,u\}|+2$
and $|M|=|M'|+1$, these contribute
$\sum_{M'}(-1)^{|M'|+1}x^{|S\setminus\{v,u\}|-2|M'|}=-\mu_{S\setminus\{v,u\}}(x)$.
\end{proof}

\subsection{The leaf expansion of a determinant}

\begin{lemma}\label{lem:det}
Let $R$ be a commutative ring, $I$ a finite index set, $N\in R^{I\times I}$, and $v\ne u$ in $I$
such that $N_{vw}=0$ and $N_{wv}=0$ for all $w\notin\{v,u\}$. Writing $N^{(v)}$ and $N^{(v,u)}$
for the principal submatrices obtained by deleting $v$, resp.\ $v$ and $u$,
\[
  \det N \;=\; N_{vv}\,\det N^{(v)}\;-\;N_{vu}N_{uv}\,\det N^{(v,u)}.
\]
\end{lemma}
\begin{proof}
Row $v$ of $N$ equals $N_{vv}e_v+N_{vu}e_u$, so by linearity of the determinant in row $v$,
$\det N=N_{vv}\det N_1+N_{vu}\det N_2$, where $N_1$ (resp.\ $N_2$) is $N$ with row $v$ replaced by
$e_v$ (resp.\ $e_u$).
\emph{(i)} In $N_1$ the row $v$ is $e_v$; after moving $v$ to the first position, $N_1$ is block
lower triangular with blocks $(1)$ and $N^{(v)}$, so $\det N_1=\det N^{(v)}$.
\emph{(ii)} In $N_2$ the column $v$ equals $N_{uv}e_u$ (its entry in row $v$ is $(e_u)_v=0$, and
the other entries of column $v$ are those of $N$). By linearity in column $v$,
$\det N_2=N_{uv}\det P$, where $P$ is $N_2$ with column $v$ replaced by $e_u$; so row $v$ and
column $v$ of $P$ both equal $e_u$. Swapping rows $u$ and $v$ of $P$ gives a matrix $Q$ with
$\det Q=-\det P$ whose row $u$ is $e_u$; expanding as in (i), $\det Q=\det Q^{(u)}$. In $Q^{(u)}$
the column $v$ is $e_v$ (its entry in row $v$ is $P_{uv}=1$, the others are $P_{wv}=0$ for
$w\ne u,v$), so expanding along column $v$, $\det Q^{(u)}=\det Q^{(u,v)}$, and the remaining
matrix $Q^{(u,v)}$ is exactly $N^{(v,u)}$ (its entries avoid rows and columns $u,v$, where $Q$
agrees with $N$). Altogether $\det N_2=N_{uv}\cdot(-\det N^{(v,u)})$, which gives the claim.
\end{proof}

\begin{lemma}\label{lem:charp}
With $S,v,u$ as in Lemma~\ref{lem:mu},
$\ \det(xI-A_S)=x\det(xI-A_{S\setminus\{v\}})-\det(xI-A_{S\setminus\{v,u\}})$.
\end{lemma}
\begin{proof}
Apply Lemma~\ref{lem:det} to $N=xI-A_S$ over $R=\Z[x]$ (or any commutative ring). Row $v$ of
$A_S$ has a single $1$, in column $u$, and $A_S$ is symmetric, so the support hypotheses hold;
$N_{vv}=x$ (there are no loops), $N_{vu}=N_{uv}=-1$, and $(xI-A_S)^{(v)}=xI-A_{S\setminus\{v\}}$,
$(xI-A_S)^{(v,u)}=xI-A_{S\setminus\{v,u\}}$.
\end{proof}

\subsection{Proof of Theorem~\ref{thm:A}}
By strong induction on $|S|$. If $G$ has no edge inside $S$ then the only matching of $G[S]$ is
$\emptyset$, so $\mu_S(x)=x^{|S|}$, while $A_S=0$ and $\det(xI-0)=x^{|S|}$. Otherwise
Lemma~\ref{lem:leaf} provides $v\in S$ with a unique neighbour $u\in S$, and Lemmas~\ref{lem:mu}
and~\ref{lem:charp} give the same recurrence for $\mu_S$ and for $\det(xI-A_S)$ in terms of the
strictly smaller sets $S\setminus\{v\}$ and $S\setminus\{v,u\}$, where the two agree by induction.
The case $S=V$ is the statement for $G$. \qed

\subsection{Proof of Theorem~\ref{thm:main}}
A tree is a forest, so $\mu(T;x)=\charp(A)$ with $A=A(T)$ real symmetric
(Theorem~\ref{thm:A}, over $R=\R$ and over $R=\C$; the matching polynomial is compatible with
change of coefficient ring). By the spectral theorem $A=U\,\mathrm{diag}(\lambda_1,\dots,\lambda_n)\,U^{*}$
with $U$ orthogonal and $\lambda_i\in\R$, hence $\charp(A)=\prod_i(x-\lambda_i)$ splits over $\R$,
has exactly $n=\deg\charp(A)$ roots with multiplicity, and its roots in $\C$ are the real numbers
$\lambda_i$: this is (a). For (b), $\sum_i\lambda_i^{\,k}=\tr\,\mathrm{diag}(\lambda_i)^k
=\tr(U\mathrm{diag}(\lambda_i)^kU^{*})=\tr A^k=\sum_v (A^k)_{vv}$, and $(A^k)_{vw}$ is the
number of walks of length $k$ from $v$ to $w$. \qed

\section{The Lean~4 formalization}

The directory \texttt{lean/} is a Lake project (toolchain \texttt{leanprover/lean4:v4.33.1},
Mathlib tag \texttt{v4.33.1}, revision \texttt{0df444a360eaa60ab8c11dca51a86af692955474}).
\begin{center}\small
\begin{tabular}{@{}ll@{}}
\toprule
File & Content\\
\midrule
\texttt{Submission/DetLeaf.lean} & Lemma~\ref{lem:det}: \texttt{TLMC3718.det\_leaf}\\
\texttt{Submission/MatchingPoly.lean} & definitions; Lemma~\ref{lem:mu}: \texttt{SimpleGraph.matchingPolyIn\_leaf}\\
\texttt{Submission/Leaf.lean} & Lemma~\ref{lem:leaf}: \texttt{SimpleGraph.IsAcyclic.exists\_leaf\_in}\\
\texttt{Submission/Charpoly.lean} & Lemma~\ref{lem:charp}, Theorem~\ref{thm:A}: \texttt{SimpleGraph.matchingPolynomial\_eq\_charpoly}\\
\texttt{Submission/Main.lean} & Theorem~\ref{thm:main}: \texttt{TLMC3718.conjecture\_00000003718}\\
\texttt{Submission/Verification.lean} & \texttt{\#print axioms} of every result\\
\bottomrule
\end{tabular}
\end{center}

\paragraph{Definitions.} For \texttt{G : SimpleGraph V} with \texttt{[Fintype V] [DecidableEq V]
[DecidableRel G.Adj]} and \texttt{s : Finset V}:
\begin{lstlisting}
def IsMatchingIn (s : Finset V) (M : Finset (Sym2 V)) : Prop :=
  M @$\subseteq$@ G.edgeFinset @$\wedge$@ (@$\forall$@ e @$\in$@ M, @$\forall$@ x, x @$\in$@ e @$\to$@ x @$\in$@ s) @$\wedge$@
    @$\forall$@ e @$\in$@ M, @$\forall$@ f @$\in$@ M, e @$\neq$@ f @$\to$@ @$\forall$@ x, x @$\in$@ e @$\to$@ x @$\notin$@ f
noncomputable def matchings (s : Finset V) : Finset (Finset (Sym2 V)) :=
  G.edgeFinset.powerset.filter (G.IsMatchingIn s)
noncomputable def numMatchings (s : Finset V) (k : @$\mathbb{N}$@) : @$\mathbb{N}$@ :=
  ((G.matchings s).filter (fun M => M.card = k)).card
noncomputable def matchingPolyIn (s : Finset V) : R[X] :=
  @$\sum$@ M @$\in$@ G.matchings s, (-1 : R[X]) ^ M.card * X ^ (s.card - 2 * M.card)
noncomputable def matchingPolynomial : R[X] := G.matchingPolyIn R Finset.univ
def adjMatrixIn (s : Finset V) : Matrix s s R :=
  (G.adjMatrix R).submatrix Subtype.val Subtype.val
\end{lstlisting}
So \texttt{G.matchingPolynomial R} is $\mu(G;x)$ with coefficients in the commutative ring $R$;
the theorem \newline \texttt{matchingPolynomial\_eq\_sum\_numMatchings} proves that it equals
$\sum_{k=0}^{n}(-1)^k\,m_k\,X^{n-2k}$ with $m_k=$\texttt{numMatchings univ k}, the classical form.
Trees are Mathlib's \texttt{SimpleGraph.IsTree} (connected and acyclic); walks are Mathlib's
\texttt{SimpleGraph.Walk}.

\paragraph{Main statements.}
\begin{lstlisting}
theorem SimpleGraph.matchingPolynomial_eq_charpoly (hG : G.IsAcyclic) :
    G.matchingPolynomial R = (G.adjMatrix R).charpoly

theorem TLMC3718.conjecture_00000003718 (T : SimpleGraph V) [DecidableRel T.Adj]
    (hT : T.IsTree) :
    (T.matchingPolynomial @$\mathbb{R}$@).Splits @$\wedge$@
    (T.matchingPolynomial @$\mathbb{R}$@).roots.card = Fintype.card V @$\wedge$@
    (@$\forall$@ z : @$\mathbb{C}$@, (T.matchingPolynomial @$\mathbb{C}$@).IsRoot z @$\to$@ z.im = 0) @$\wedge$@
    (@$\forall$@ k : @$\mathbb{N}$@, ((T.matchingPolynomial @$\mathbb{R}$@).roots.map (fun r => r ^ k)).sum =
      @$\sum$@ v, (Fintype.card {p : T.Walk v v | p.length = k} : @$\mathbb{R}$@))
\end{lstlisting}
Here \texttt{Polynomial.Splits} is Mathlib's predicate ``product of a constant and linear
factors''; the four conjuncts are Theorem~\ref{thm:main}(a) (three forms) and (b). The individual
results are \texttt{matchingPolynomial\_splits}, \texttt{card\_roots\_matchingPolynomial},
\texttt{exists\_real\_of\_isRoot} / \texttt{im\_eq\_zero\_of\_isRoot} and
\texttt{sum\_roots\_pow\_eq\_card\_closedWalks}.

\paragraph{Audit.} \texttt{lake build} succeeds; the sources contain no \texttt{sorry},
\texttt{native\_decide}, \texttt{axiom}, \texttt{unsafe}, \texttt{opaque},
\texttt{implemented\_by} or \texttt{extern}. \texttt{Submission/Verification.lean} prints, for
every theorem listed above,
\begin{lstlisting}
'TLMC3718.conjecture_00000003718' depends on axioms: [propext, Classical.choice, Quot.sound]
\end{lstlisting}

\paragraph{Reproduction.} In \texttt{lean/}: \ \texttt{lake exe cache get \&\& lake build}.
This report compiles with \texttt{tectonic report.tex} (or \texttt{pdflatex}).

\section*{References}
\begin{itemize}\small
\item O.~J. Heilmann, E.~H. Lieb, \emph{Theory of monomer--dimer systems}, Comm.\ Math.\ Phys.\ 25 (1972), 190--232.
\item C.~D. Godsil, I.~Gutman, \emph{On the theory of the matching polynomial}, J.\ Graph Theory 5 (1981), 137--144.
\item C.~D. Godsil, \emph{Algebraic Combinatorics}, Chapman \& Hall, 1993, Ch.~1 and~6.
\end{itemize}

\end{document}
